\section{Conclusion}
\label{sec:conclusion}

\subsection{Summary}

We have presented Multi-scale MCMC Option B for redistricting, implemented
in \texttt{redist-ensemble} and exposed via \texttt{--search multiscale}
in \texttt{redist-cli}.
The key results are:

\begin{enumerate}

  \item \textbf{County adjacency from tract graph (exact, no extra data).}
  The county adjacency graph required for coarse \recom{} moves is derived
  exactly from the existing tract graph by scanning cross-county tract edges.
  Two counties are adjacent in the derived graph if and only if they are
  geographically adjacent.
  No additional data files, shapefiles, or census downloads are required.
  The derivation runs in $O(|E|)$ time ($<1$ ms for Texas).

  \item \textbf{Mixing speed improvement for large $k$.}
  On Texas ($k = 38$, $T = 2000$ steps), \ms{} reduces lag-100
  autocorrelation of the edge-cut statistic from $0.51$ to $0.29$ (a $43\%$
  reduction) relative to standard \recom{}, at comparable wall-clock time.
  This yields a $\approx 1.8\times$ increase in effective sample size per
  unit time.

  \item \textbf{Three resolution options via GEOID prefix matching.}
  The resolution system defines Options A (block group $\to$ tract), B
  (tract $\to$ county, implemented), and C (block group $\to$ county), all
  sharing the \texttt{derive\_partition} helper.
  Options A and C require block-group adjacency data not currently loaded
  by the pipeline and are defined for future implementation.

\end{enumerate}

\subsection{When to Use Multi-scale MCMC}

\ms{} is the recommended ensemble method when:
\begin{itemize}
  \item $k \geq 30$ (the large-$k$ mixing problem is acute); or
  \item The state has a strong county-district correspondence (many states
        have constitutional or statutory requirements to respect county
        boundaries); or
  \item Ensemble analysis requires effective independence between samples
        (litigation support, academic publication, legislative compliance).
\end{itemize}

Standard \recom{} remains appropriate for:
\begin{itemize}
  \item $k \leq 14$ (NC, WI, OR, etc.) where single-scale mixing is adequate;
  \item Rapid exploratory runs where ESS precision is not critical.
\end{itemize}

\subsection{Parameter Guidelines}

Based on the $\alpha$-sweep results in Section~\ref{sec:mixing}:
\begin{itemize}
  \item $\alpha = 0.3$ is the recommended default for states with $k = 30$--$45$.
  \item $\alpha = 0.35$ may improve mixing for $k \geq 50$ (CA, $k = 52$).
  \item $\alpha < 0.2$ provides minimal benefit over standard \recom{}.
  \item $\alpha > 0.45$ increases rebalance failure rates and degrades
        effective acceptance.
\end{itemize}

\subsection{Future Directions}

Three natural extensions of the current implementation:

\paragraph{Multi-scale Forest ReCom.}
Replacing the fine \recom{} chain with \fr{}~\citep{dellaLibera2026forestRecom}
would provide distributional guarantees for the fine-level component.
The Metropolis-Hastings acceptance ratio used in \fr{} is well-defined on
the fine plan space; combining it with coarse moves requires care in
defining the joint acceptance criterion.

\paragraph{Options A and C.}
Implementing block-group adjacency loading would enable Options A and C,
providing finer spatial resolution for states with heterogeneous tract
populations (e.g.\ California, where coastal tracts are very small).

\paragraph{Adaptive $\alpha$.}
The coarsening parameter $\alpha$ could be adapted during the run based on
the rebalance failure rate: decrease $\alpha$ when failures exceed a threshold,
increase when failures are rare.
This would automate the $\alpha$ selection currently done manually.

\subsection{Implementation Reference}

The complete implementation is in \texttt{crates/redist-ensemble/src/multiscale.rs}.
The CLI integration is in \texttt{crates/redist-cli/src/compositor/multiscale.rs}.
The GEOID-prefix coarsening helper is in
\texttt{crates/redist-core/src/adjacency\_loader.rs}
(\texttt{derive\_partition} and \texttt{build\_county\_coarsening}).

All runs are reproducible from the manifest entry \texttt{plan\_resolution: ``tract-county''}
and the base seed $b$, using the deterministic step-level seed schedule
(Section~\ref{sec:algorithm}).
